\documentclass[11pt,a4paper]{article}
\usepackage[margin=25mm]{geometry}
\usepackage{amsmath,amssymb,amsthm,hyperref}
\setlength{\emergencystretch}{2em}
\newtheorem{lemma}{Lemma}
\newtheorem{theorem}[lemma]{Theorem}
\newcommand{\gen}[1]{\langle#1\rangle}
\newcommand{\Alt}{\operatorname{Alt}}
\newcommand{\Sym}{\operatorname{Sym}}
\title{The dihedral head cannot meet a natural alternating tail}
\author{Human Compression project: bounded L044 proof component}
\date{}
\begin{document}\maketitle
All permutations act on the right; $g^h=h^{-1}gh$.
A subgroup described as \emph{pure} on a set fixes its complement
pointwise. Products and intersections below are of actual subgroups,
not of abstract isomorphism types or unproved projected C-markings.

\begin{lemma}[A seven-point alternating subgroup]
Let $X$ be disjoint from $\{1,2,3,4,5\}$ and contain $6,y$,
with $y\ne6$. Put
\[
 l=(1\ 2)(3\ 4),\qquad s=(2\ 5)(1\ 3),\qquad
 t=(1\ 3)(2\ 4)(5\ 6)\delta,
\]
where $\delta$ fixes $1,\ldots,6$ and is supported on $X$.
Let $e=(6\ y)u$ be an involution supported on $X$, where $u$
fixes $6,y$. Then $\gen{l,s,t,e}$ contains the pure alternating
group on $Z=\{1,2,3,4,5,6,y\}$.
\end{lemma}
\begin{proof}
The subgroup $D=\gen{l,s,s^t}$ is pure on $Z_0=\{1,\ldots,6\}$.
Indeed $s^t=(4\ 6)(1\ 3)$. The group $\gen{l,s}$ is
transitive on $\{1,\ldots,5\}$: its product is a $5$-cycle.
Since $s^t$ moves $6$, $D$ is transitive on $Z_0$.
Its stabilizer of $6$ contains $\gen{l,s}$ and is transitive
on the other five points. Thus $D$ is $2$-transitive on $Z_0$.

Put $a=s^t$. One has $a^e=(4\ y)(1\ 3)$, so $aa^e$ is
a pure $3$-cycle on $\{4,6,y\}$. Conjugating it by $D$, which
fixes $y$, gives the $3$-cycles on $\{v,w,y\}$ for all
ordered distinct $v,w\in Z_0$. These generate $\Alt(Z)$:
they include every $3$-cycle involving $y$, and products of
two such cycles express those not involving $y$.
All these operations occur inside $\gen{l,s,t,e}$.
\end{proof}

\begin{theorem}[The remaining decorated dihedral case is impossible]
Let $G=\gen{\rho_0,\ldots,\rho_{r-1}}\leq\Alt(\Omega)$
be a string C-group, where
$\Omega=\{1,2,3,4,5\}\sqcup X$ and $|X|\geq5$.
Suppose its actual first three generators are
\[
 \rho_0=l=(1\ 2)(3\ 4),\qquad
 \rho_1=s=(2\ 5)(1\ 3),\qquad
 \rho_2=t=\lambda(5\ 6)\delta,
 \quad \lambda=(1\ 3)(2\ 4),
\]
where $6\in X$, and $\delta$ is an involution supported on
$X\setminus\{6\}$. Suppose
$K=\gen{\rho_3,\ldots,\rho_{r-1}}$ is the full pure natural
alternating group $\Alt(X)$. No such marking exists.
\end{theorem}
\begin{proof}
Since $\lambda$ is even and $t$ is even, $\delta$ is odd.
Every $\rho_j$ with $j\geq4$ fixes $5$ and commutes with $t$,
so fixes $6=5^t$ and centralizes $\delta$. Transitivity of $K$
therefore forces $e=\rho_3$ to move $6$. Write
$e=(6\ y)u$, where $y\in X\setminus\{6\}$ and $u$ is an
odd involution on $X\setminus\{6,y\}$.

We first identify an actual subgroup used in the intersection
argument. Set $B=\{5\}\cup X$ and $R=\gen{t,K}$.
Then $R\leq\gen\lambda\times\Alt(B)$. Its projection to $B$
is the full $\Alt(B)$: it contains $\Alt(X)$ and its conjugate
under $(5\ 6)\delta$, namely the alternating group on
$B\setminus\{6\}$. These two natural alternating groups,
on sets of size at least five with all but one point in common,
generate $\Alt(B)$ by their $3$-cycles. The other projection
is onto $\gen\lambda$. If the projection to $\Alt(B)$ were
injective, it would give a nontrivial homomorphism from the
perfect group $\Alt(B)$ to $C_2$. This is impossible.
Its kernel is consequently $\gen\lambda$, and
\[
                         R=\gen\lambda\times\Alt(B).
\]

Let $P=\gen{l,s,t,e}$ and $Z=\{1,\ldots,6,y\}$.
By Lemma~1, $P$ contains pure $\Alt(Z)$. In particular it
contains both $\lambda$ and the pure $3$-cycle $c=(5\ 6\ y)$.
Both elements also lie in $R$. Full ambient IP gives
\[
                           P\cap R=\gen{t,e}.
\]
This is a finite dihedral group, in which all elements of odd
order belong to the cyclic rotation subgroup and commute.
If $\delta$ moved $y$, the two $3$-cycles $c,c^t$ would have
supports $\{5,6,y\}$ and $\{5,6,y^\delta\}$, respectively,
with exactly two common points. They do not commute.
Both belong to $\gen{t,e}$, a contradiction. Hence
$y^\delta=y$.

It follows that $P$ preserves $Z$. Any element in the kernel
of its action on $Z$ is supported on $X\setminus\{6,y\}$
and is even, since $G\leq\Alt(\Omega)$. It therefore belongs
to $K=\Alt(X)$. But full ambient IP also gives
\[
                              P\cap K=\gen e.
\]
The nonidentity element $e$ moves $6$ and $y$, so this kernel
is trivial. Thus the action of $P$ on $Z$ is faithful.
It contains pure $\Alt(Z)$, and $t$ induces the odd permutation
$\lambda(5\ 6)$, so its image is $\Sym(Z)$.

In particular every element of $P$ inducing an even
permutation on $Z$ is the unique pure alternating permutation
with that action. Apply this to $te$: both $t$ and $e$ are odd
on $Z$, so $te$ acts trivially off $Z$. Their exterior parts
are $\delta$ and $u$, respectively, whence $\delta u=1$ and
$u=\delta$. Consequently $e=(6\ y)\delta$ centralizes
$\delta$, which already centralizes every other generator
of $K$. The permutation $\delta$ centralizes the transitive
group $K$ and fixes $6$, so it fixes all of $X$.
This contradicts its oddness and proves the theorem.
\end{proof}

\paragraph{Source application and limits.}
The source is Cameron--Fernandes--Leemans--Mixer,
\emph{Highest rank of a polytope for $A_n$},
arXiv:1605.09173v1. This fills the remaining natural-alternating-tail
case in Proposition~5.14 (physical pages~28--29) and the
$k=2$ residual isolated in the accepted replacement of the
intransitive-side equality step of Proposition~5.26
(physical pages~33--34). The two descriptions have the same
literal marking after relabelling. The dihedral head, the
source's attachment strategy and the full-IP requirement are
inherited; the seven-point alternating-subgroup and faithful
kernel argument above are the present replacement proof.

The six-point projection of $\gen{l,s,t}$ is the decorated
$S_5$ case identified in the earlier local audit; it must not
be silently replaced by either $A_5$ row of source Table~2.
Our proof does not need that structural identification.
It uses no smaller-degree rank induction, no projected IP,
no finite census and no new formal theorem. It does not
by itself settle the other source branches or confer whole
L044 recovery or question credit.
\end{document}
